\section{Conclusion}
\begin{frame}{Summary (1/2)}
\begin{itemize}
\item Proximal gradient methods with spectral steps (NSPG, NSPGH) often converge to better local minima than coordinate descent (PGCCD) despite weaker theoretical guarantees;
\item The combination NSPG+PGCCD as an inner solver within the CPSI($k$) framework yields the best support recovery and parameter estimation, especially under constant correlation;
\item Cross-validation strategy matters: inverse paths benefit NSPG; adaptive paths are best for PGCCD and combined methods.
\end{itemize}
\end{frame}

\begin{frame}{Summary (2/2)}
\begin{itemize}
\item The variable metric variant (VMNSPG) converges quickly but occasionally traps in poor minima;
\item Compared to L0Learn, the new methods provide more robust solutions in difficult correlation settings with competitive computational efficiency;
\item Tridiagonal masking reveals the importance of problem topology: dense global correlations are harder, but removing them can sometimes harm fine differentiation.
\end{itemize}
\end{frame}

\begin{frame}{Future Work}
\begin{itemize}
\item Extend to the $\ell_0\ell_1$ and $\ell_0\ell_2$;
\item Investigate structured (e.g., tridiagonal) variable metric proximal schemes to capture more local second-order information without full Hessian. Each proximal subproblem must be solved iteratively;
\item Investigate scenarios with partially sparse true Hessians;
\item Rigorous statistical theory for the quality of the solutions delivered by the proposed algorithms (is it even possible?).
\end{itemize}
\end{frame}